\section{The Cost Frontier}
\sectionkicker{COST EFFICIENCY}
\begin{wideflow}
{\Large\bfseries 値段の下がった最前線――SolとLuna、そしてプロンプトキャッシュの改善}\par
\smallskip
{\color{SurveyMuted}9月22日のSol/Luna公開とキャッシュ改善を、半額の料金表と提供範囲、評価の留保と切り分けて示す。}\par
\end{wideflow}

9月22日、OpenAIはGPT-6の系列にSolとLunaを加えた\cite{gpt6solluna}。Astraと同じ流儀で鍛え、その高い性能を速く安い二つのモデルで届ける内容で、最上位モデルの交代ではない。APIではgpt-6-solとgpt-6-lunaの名で呼ばれ、ChatGPT WorkとCodexでは有料の各プランから同日に入り、Lunaは無料・Goプランの利用者にもデスクトップアプリで利用できる。一方でChat（対話アプリ）への広がりはその日のうちの順次とされ、見えない利用者もいた。料金は5.6世代の販促価格から半額となり、Solは入力100万トークンあたり4ドルが2ドル、出力20ドルが10ドル、Lunaは入力0.20ドルが0.10ドル、出力1.20ドルが0.50ドルと示された。請求の実測は本号では確かめていない。

ベンチマーク結果は発表側の計測による。業務の流れをこなすAutomationBenchではSolのxhigh設定が33.2\%を刻み、1タスクあたり27セントと置かれる。相手のOpus 5の26.9\%を、費用では11分の1ほどで上回るとされる。長時間タスク向けのAgents' Last Examでは56.4\%、コーディングベンチマークのFrontierCodeでは前世代のGPT-5.6 Solを離してFable 5.1の高effort設定に並び、DeepSWEでは68.8\%でFable 5の69.9\%に1.1差まで迫り、費用は8割ほど軽いとされる。実在のやり取りで誤りを数える社内の評価では、Solの誤りが前世代（GPT-5.6 Sol）の半分ほどに減り、Astra並みの確かさに近づくとされる。いずれも発表側の計測方法であり、相手の数は公開の報告からの引用である。

同じ日に、長期間動作するエージェントのためのプロンプトキャッシュの改良も出た\cite{gpt6caching}。使い回す共通プロンプト先頭の一致で費用と待ち時間を抑える仕組みが強まり、30分間の再利用窓での再利用に割引がつく。どれだけ効いているかを見るダッシュボード、外した要求を説明するキャッシュミス診断ツール、キャッシュ境界を利用者が指定する手段、途中でreasoning effortを変えてもキャッシュを保つ手立てが添えられた。GitHubでは数十億のリクエストを通じて新たに作るトークンの割合が半分以上減ったとされるが、これも発表側の引用である。公開の場では値下げを喜ぶ声と、長時間タスクでのもたつきを訴える声が交じり、後者は一次発表では確かめも否定もできない。

脇には、前号から持ち越したDeepSeekの話に区切りがついた。公式の料金表にdeepseek-flash（V4.1-Flash相当）とdeepseek-v4-pro（V4-Pro-0813相当）が並び、旧モデル名へのAPIリクエストは新しいFlashに回してFlashの値で数えると明記された\cite{deepseekdocs}。前号で開いたままだった一次資料の穴はこの一枚で塞がり、V4.1-Proというモデルの登場や重みの退役は書かれていないので本号でも言わない。正確な切替時刻はこの公式ページからは読めない。

\begin{claimboundary}[資料と評価の境界]
料金表は発表時点の公式ページであり、請求の実測ではない。評価の詳細は本号の検証範囲外であり、数値はベンダー側の報告として独立した再現を待つ。相手の数値は公開報告からの引用であり、横断的な順位づけには使わない。音声まわりの提供の細目はこの発表の範囲外である。
\end{claimboundary}

\begin{communitynote}[X / Community observation --- context only]
公開投稿では、値下げを喜ぶ声と、長時間タスクでのもたつきを訴える独立した報告\cite{w39x-c1-regression}が交じる。うかがえるのは価格への関心とエージェント運用の摩擦であり、仕様やベンチマーク、料金、利用条件の根拠にはならない。締め切り後の後追いの報告は背景情報であり、通常の集計には含めない。
\end{communitynote}
